%=====================================================================
%  Taxonomy figure (signature figure, place right after Introduction).
%  Built around the THESIS: shared-with-single-agent dims vs MARL-specific dims.
%  Uses forest package (\usepackage[edges]{forest} in main.tex).
%  This is a DRAFT layout -- tune colors/spacing at polish time.
%=====================================================================
\begin{figure*}[t]
\centering
\footnotesize
\begin{forest}
  for tree={
    grow=east, parent anchor=east, child anchor=west,
    draw, rounded corners, align=left, edge={-latex},
    l sep=10mm, s sep=2mm, anchor=west,
  }
  [Robust\\MARL, fill=gray!15
    [\textbf{I. Perturbation source}\\(shared w/ single-agent RL), fill=blue!8
      [{Environment / model uncertainty (\S3)\\\,DRMG, sample complexity, minimax}]
      [State / observation perturbation (\S4)]
      [Action perturbation (\S5)]
      [Reward poisoning (\S5)]
    ]
    [{\textbf{II. MARL-specific dimensions}\\(absent in single-agent RL; \emph{core of this survey})}, fill=red!8
      [Communication attacks \& noise (\S6)]
      [Untrustworthy teammates /\\Byzantine / fault tolerance (\S7)]
      [{Curse of multiagency\\(theory, \S3)}]
    ]
    [\textbf{III. Emerging frontiers}, fill=green!8
      [LLM multi-agent safety (\S8)]
      [Offline / distribution shift (\S9)]
    ]
    [\textbf{IV. Cross-cutting}, fill=yellow!12
      [Benchmarks \& evaluation (\S10)]
      [Applications (\S11)]
    ]
  ]
\end{forest}
\caption{Taxonomy of robust MARL. Branch~I collects perturbation dimensions
inherited from single-agent robust RL; Branch~II collects dimensions that are
\emph{specific to the multi-agent setting} and constitute the core argument of
this survey. A secondary methodological axis (distributionally robust
optimization, adversarial training, certified robustness, regularization) cuts
across all branches and is indicated per work in the section tables.}
\label{fig:taxonomy}
\end{figure*}
